\begin{abstract}
本文提出并形式化一种“扫描--投影读出”框架，用以描述在有限分辨率观测下，“存在/不存在”二值命题如何作为同一本体状态的读出投影而出现，并在黄金分支（golden branch）上获得可计算的稳定结构。核心构造为：以单位化全局态为本体载体，引入非周期的酉扫描算子生成访问序；以有限窗口投影将连续相位轨道读出为二字母机械词；在语法约束（无相邻“11”）下得到黄金均值移位（golden mean shift）稳定子空间，其维数随窗口长度按斐波那契律增长，诱导出一类“黄金存在投影”与“黄金存在态”。本文进一步提出三通道稳定性分解：$\varph$ 通道对应语法合法性与自相似增长；$\pi$ 通道对应周期闭合与轨道计数结构；$e$ 通道对应解析稳定域与对数读出代价。由此可将“只能观测到存在态”的命题严格表述为对稳定投影条件化后的统计学陈述。作为可检验的计算层结果，本文给出基于 Zeckendorf 表示与窗口截断的折叠映射 $\mathrm{Fold}_m$，将 $2^m$ 微观态压缩至 $\card{X_m}$ 个稳定类型，并在 $m=6$ 的最小例中得到 $64\to 21$ 的刚性压缩与边界/循环 $3\oplus 18$ 分裂。最后，本文在算术骨架层面讨论素数作为乘法生成元、$\pi$ 作为闭合相位几何不变量、$e$ 作为乘法权重到加法代价的唯一连续同态基，并给出可复现的算法接口与可证伪的分层预测清单。
\end{abstract}

\begin{sloppypar}
\textbf{关键词：} 扫描--投影读出；黄金存在态；黄金均值移位；Sturmian 词；Zeckendorf 表示；折叠压缩；周期闭合；对数读出代价；素数谱
\end{sloppypar}

\tableofcontents
\newpage

